\begin{center}
\textbf{DATOS:} \quad \textbf{FORMULA:}
\end{center}

\[ \lambda_1 = 2 \, \text{R/h} \quad  \lambda_2 = \left( \frac{d_1}{d_2} \right)^2 \times \lambda_1 \]

\[ d_1 = 1 \, \text{m} \]

\[ \lambda_2 = \, ? \]

\[ \lambda_2 = \frac{1}{5} \times 2 = 0.08 \, \text{R/h} = 80 \, \text{mR/h} \]

%Ahora calcularemos el tiempo que tardaríamos en recibir 100 mrem.
\textbf{Ahora calcularemos el tiempo que tardaríamos en recibir 100 mrem.}

\[
\begin{aligned}
80 \, \text{mrem} & = 1 \, \text{hr} \\
100 \, \text{mrem} & = \frac{1 \times 100}{80} = 1.25 \, \text{hr} = 1 \, \text{hr} \, 15 \, \text{min}
\end{aligned}
\]

%Si en lugar de trabajar 1 hr 15 min trabajamos sólo 40 min, ¿qué dosis recibí?
\textbf{Si en lugar de trabajar 1 hr 15 min trabajamos solo 40 min, ¿qué dosis recibió?}

\[
X = \frac{80 \times 40}{60} = 53.3 \, \text{mrem}
\]

%La constante GAMMA de Desintegración o emisividad del Ir-192 es:
\textbf{La constante GAMMA de desintegración o emisividad del Ir-192 es:}

\[
0.55 \, \text{R/h a 1 metro por cada curie}
\]

\begin{center}
\textbf{UNA FUENTE CON} \quad \textbf{NOS DA A UN METRO}

\begin{tabular}{ccc}
2 Ci & 0.55 x 2 = 1.1 R/h\\
10 Ci & 0.55 x 10 = 5.5 R/h\\
35 Ci & 0.55 x 35 = 19.25 R/h\\
100 Ci & 0.55 x 100 = 55 R/h\\
\end{tabular}
\end{center}

¿Cuánto espesor de plomo se necesita para evitar que una fuente de Ir-192 de 53.3 Ci de más de 100 mR/h a un metro?

%As you requested, the LaTeX code is given without additional embellishments, strictly representing the image content with corrected context where necessary.
